\documentclass[../../../include/open-logic-section]{subfiles}

\begin{document}

\olfileid[id]{sth}{spine}{rank}
\olsection{Peringkat}

Sekarang kita telah mendefinisikan tahap-tahap sebagai $V_\alpha$, dan kita
tahu bahwa setiap himpunan merupakan himpunan bagian dari suatu tahap, kita
dapat mendefinisikan \emph{peringkat} suatu himpunan. Secara intuitif,
peringkat $A$ adalah tahap pertama tempat $A$ terbentuk. Lebih tepatnya:

\begin{defn}\ollabel{defnsetrank}
Untuk setiap himpunan $A$, $\setrank{A}$ merupakan ordinal terkecil $\alpha$
sedemikian sehingga $A \subseteq V_\alpha$.
\end{defn}
\begin{prop}\ollabel{ranksexist}
	$\setrank{A}$ ada, untuk sebarang $A$.
\end{prop}
\begin{proof}
	Ditinggalkan sebagai latihan.
\end{proof}
\begin{prob}
	Buktikan \olref[sth][spine][rank]{ranksexist}.
\end{prob}
\noindent
Pengurutan baik peringkat memungkinkan kita membuktikan beberapa hasil
penting:
\begin{prop}\ollabel{valphalowerrank}
Untuk setiap ordinal $\alpha$,
$V_\alpha = \Setabs{x}{\setrank{x} \in \alpha}$.
\end{prop}

\begin{proof}
Jika $\setrank{x} \in \alpha$, maka
$x \subseteq V_{\setrank{x}} \in V_\alpha$, sehingga $x \in V_\alpha$
karena $V_\alpha$ bersifat poten (dengan menggunakan
\olref[Valphabasic]{Valphabasicprops} beberapa kali). Sebaliknya, jika
$x \in V_\alpha$, maka $x \subseteq V_\alpha$, sehingga
$\setrank{x} \leq \alpha$; sekarang induksi transfinit sederhana menunjukkan
bahwa $x \notin V_{\setrank{x}}$.
\end{proof}
\begin{prob}
	Lengkapi induksi transfinit sederhana yang disebutkan dalam
	\olref[sth][spine][rank]{valphalowerrank}.
\end{prob}

\begin{prop}\ollabel{rankmemberslower}
Jika $B \in A$, maka $\setrank{B} \in \setrank{A}$.
\end{prop}

\begin{proof}
$A \subseteq V_{\setrank{A}} = \Setabs{x}{\setrank{x} \in
\setrank{A}}$ berdasarkan \olref{valphalowerrank}.
\end{proof}
\noindent
Dengan menggunakan fakta ini, kita dapat menetapkan suatu hasil yang
memungkinkan kita membuktikan berbagai hal tentang \emph{semua himpunan}
melalui suatu bentuk induksi:

\begin{thm}[Skema Induksi-$\in$]
Untuk sebarang formula $\phi$:
\[
	\forall A((\forall x \in A)\phi(x) \lif \phi(A)) \lif \forall A \phi(A).
\]
\end{thm}

\begin{proof}
Kita akan membuktikan kontraposisinya. Andaikan
$\lnot \forall A \phi(A)$. Berdasarkan Induksi Transfinit
(\olref[ordinals][basic]{ordinductionschema}), terdapat suatu objek yang
tidak bersifat $\phi$ dan peringkatnya sekecil mungkin; yakni suatu $A$ sedemikian
sehingga $\lnot \phi(A)$ dan
$\forall x(\setrank{x} \in \setrank{A} \lif \phi(x))$. Sekarang, jika
$x \in A$, maka $\setrank{x} \in \setrank{A}$ berdasarkan
\olref{rankmemberslower}, sehingga $\phi(x)$; yakni
$(\forall x \in A)\phi(x) \land \lnot \phi(A)$.
\end{proof}\noindent
Berikut cara informal untuk menjelaskan hasil yang kuat ini. Katakan bahwa
$\phi$ bersifat \emph{herediter} jika dan hanya jika, setiap kali semua
!!{element} suatu himpunan bersifat $\phi$, himpunan itu sendiri juga
bersifat $\phi$. Maka Induksi-$\in$ memberi tahu Anda hal berikut: jika
$\phi$ bersifat herediter, setiap himpunan bersifat $\phi$.

Untuk mengakhiri pembahasan peringkat (untuk saat ini), kita akan membuktikan
beberapa klaim yang telah beberapa kali kita singgung sebelumnya.

\begin{prop}\ollabel{ranksupstrict}
$\setrank{A} = \supstrict_{x \in A}\setrank{x}$.
\end{prop}

\begin{proof}
Misalkan $\alpha = \supstrict_{x \in A}\setrank{x}$. Berdasarkan
\olref{rankmemberslower}, $\alpha \leq \setrank{A}$. Namun, jika $x \in A$,
maka $\setrank{x} \in \alpha$, sehingga $x \in V_\alpha$ berdasarkan
\olref{valphalowerrank}, dan karenanya $A \subseteq V_\alpha$, yakni
$\setrank{A} \leq \alpha$. Maka $\setrank{A} = \alpha$.
\end{proof}

\begin{cor}\ollabel{ordsetrankalpha}
Untuk setiap ordinal $\alpha$, $\setrank{\alpha} = \alpha$.
\end{cor}

\begin{proof}
Untuk induksi transfinit, andaikan $\setrank{\beta} = \beta$ bagi semua
$\beta \in \alpha$. Sekarang
$\setrank{\alpha} = \supstrict_{\beta \in \alpha}\setrank{\beta} =
\supstrict_{\beta \in \alpha}\beta = \alpha$ berdasarkan
\olref{ranksupstrict}.
%	First note that $\setrank{\alpha} \neq \beta$ for any $\beta \in
%	\alpha$. For otherwise we would have $\beta = \setrank{\beta} \in
%	\setrank{\alpha} = \beta$ by \olref{rankmemberslower}, i.e.,
%	$\beta \in \beta$, a contradiction.
%
%	Now note that $\alpha \subseteq V_\alpha$. For $\alpha =
%	\Setabs{\beta}{\beta \in \alpha}$ by
%	\olref[sfr][ordinals][basic]{ordissetofsmallerord}. And if
%	$\beta \in \alpha$ then $\beta \subseteq V_{\setrank{\beta}} =
%	V_\beta \in V_\alpha$, hence $\beta \in V_\alpha$, by
%	\olref[sfr][spine][Valphabasic]{Valphabasicprops} twice.
\end{proof}

Terakhir, berikut bukti singkat bagi hasil yang dijanjikan pada akhir
\olref[foundation]{sec}, bahwa $\ZFminus$ membuktikan implikasi
$\emph{Regularitas} \Rightarrow \emph{Fondasi}$. (Perhatikan bahwa gagasan
``peringkat'' dan \olref{rankmemberslower} tersedia untuk digunakan dalam
bukti ini karena---sebagaimana disebutkan pada awal bagian ini---keduanya
dapat disajikan dengan menggunakan $\ZFminus + \text{Regularitas}$.)

\begin{prop}[dengan bekerja dalam $\ZFminus + \text{Regularitas}$]\ollabel{zfminusregularityfoundation} Fondasi berlaku.
\end{prop}

\begin{proof}
Tetapkan $A \neq \emptyset$, dan suatu $B \in A$ yang peringkatnya sekecil
mungkin. Jika $c \in B$, maka $\setrank{c} \in \setrank{B}$ berdasarkan
\olref{rankmemberslower}, sehingga $c \notin A$ berdasarkan pemilihan $B$.
\end{proof}

\end{document}
